
\subsubsection*{Erreichter Stand}

Der in Abschnitt~\ref{lösungsweg-algorithmik} beschriebene Navigationsaufbau (A*, D* Lite, DWA, Pose\-Estimator) ist in der Simulation funktional vollständig und gegen die Testfälle \texttt{ST-SIM-001} bis \texttt{ST-SIM-006} verifiziert, ebenso die darüberliegende Sicherheitsebene (\texttt{SafetyMonitor}, \texttt{FaultMonitor}). Parallel dazu wurde die in der Softwaredokumentation zunächst nur als Vorschlag skizzierte Missionsablaufsteuerung (Tisch einlernen, Untersetzer ausgeben/einsammeln, Not-Halt) in einem eigenen Simulationsstand (\texttt{coasterbot\_V\_01\_00}) umgesetzt und erreicht dort inklusive Hindernisumfahrung eine Positioniergenauigkeit von rund \SI{1,5}{\centi\meter}. Auf der realen Hardware läuft bereits ein eigenständiger, direkt gegen die reale HAL arbeitender Punkt-zu-Punkt-Navigator (\texttt{Navigation}), der die Skid-Steer-Einschränkung (nur Geradeausfahrt und Drehung auf der Stelle, siehe Abschnitt~\ref{lösungsweg-elektronik}) berücksichtigt. Eine direkte Übernahme des simulierten A*/D*-Lite/DWA-Stapels auf den ressourcenbeschränkten Pico~2 steht noch aus.

\subsubsection*{Nutzen und Weiterentwicklung}

Bevor die algorithmischen Entscheidungen auf der realen Hardware ausgeführt wurden, konnten diese simulationsgestützt verifiziert werden.
Dies ist sichtbar am Beispiel der Odometrie: Die erste reale Iteration der Firmware verzichtete, wie ursprünglich geplant, auf Radencoder und integrierte die Position allein aus dem Beschleunigungssensor. Im Hardwaretest erwies sich die so erreichbare Genauigkeit als unzureichend (akkumulierter Fehler im Dezimeterbereich). Da der in der Simulation bereits erprobte \texttt{PoseEstimator} (Komplementärfilter aus Radodometrie und Gyroskop, verifiziert auf wenige Zentimeter bis Millimeter Abweichung bei Geradeausfahrt) zeigte, dass eine encoderbasierte Lösung zuverlässig funktioniert. Deshalb wurden nachträglich Gabellichtschranken an zwei Motoren ergänzt (siehe Abschnitt~\ref{lösungsweg-elektronik}). Hierbei hat ein simulationsinterner Ansatz direkt eine Hardwareänderung begründet.

Den physischen Rahmen dafür liefern die Ergebnisse der Elektronik- und Mechanikentwicklung (Abschnitte~\ref{lösungsweg-mechink} und~\ref{lösungsweg-elektronik}): Ohne die aufgebaute Elektronik (Recheneinheit, Motortreiber, Sensorik und Energieversorgung) als reale Plattform zur Verifikation, wären die nötigen Änderungen an der Odometrie, Navigation oder Fehlererkennung nicht sichtbar geworden. Besonderen Nutzen für die Weiterentwicklung bringt dabei die eigens konstruierte Untersetzer-Ausgabe-/Aufnahmemechanik: Statt eines der generisch betrachteten Prinzipien (Sauger, Greifer, manuelle Einlage) wurde eine kombinierte Lösung aus magnetbestückter Hubscheibe und Schubkurbeltrieb entwickelt, ergänzt um einen 90°-indexiert drehenden Ausschieber, der den Untersetzer beim Ablegen ausgerichtet auf einen Führungssteg schiebt und beim Aufnehmen den Weg für den Heber freigibt. Diese Mechanik wird über einen einzigen Servomotor angesteuert und besitzt damit bewusst nur einen Freiheitsgrad. Deshalb genügt zur Steuerung eine einfache servogesteuerte Ablaufprozedur, und der in der Alternativenbeschreibung noch mit inverser Kinematik und RRT* skizzierte, deutlich aufwendigere Ansatz (Abschnitt~\ref{lösungsweg-algorithmik}) war für die gewählte mechanische Lösung nie erforderlich.

Bei der Missionsablaufsteuerung bildet die Simulation inzwischen nicht mehr jeden Aspekt ab, der bereits Teil der realen Software ist: Im hier dokumentierten Hauptsimulationsstand existiert sie nur als Vorschlag, während sie im parallelen Stand \texttt{coasterbot\_V\_01\_00} sowie in der Firmware bereits aktiv weiterentwickelt wird.

Für die Weiterentwicklung des Projekts ergeben sich nach Möglichkeit daraus folgende nächste Schritte: die Zusammenführung der beiden Simulationsstände zu einem gemeinsamen Navigations- und Missionsmodell, sowie die schrittweise Übertragung des verifizierten Navigationssaufbaus auf die reale Firmware Zunächst kann der bereits erfolgreich erprobte \texttt{PoseEstimator}-Ansatz, danach die D*-Lite-/DWA-basierte Hindernisvermeidung, eingesetzt werden insofern ausreichend Rechenzeit auf dem Pico~2 zur Verfügung steht.
